
We investigated entropy-guided selective SWA conversion in Llama-2-7B as a zero-shot, no-fine-tuning approach to reducing attention cost. The work proceeded in two stages.

\textbf{Stage 1 (h-e1 --- Prerequisite Validation):} Per-layer attention entropy, computed via head-mean pooling over a 100-sequence WikiText-103 calibration set, reveals that Llama-2-7B layers exhibit strong attention concentration (Gini = 0.6829, top-10\% token share = 71.72\%) uniformly across 200 evaluation examples. Head-mean pooling captures this concentration 32\% more strongly than head-max pooling (Gini 0.681 vs. 0.466). Entropy rankings are stable across calibration subsets (Spearman $\rho \geq$ 0.8 confirmed), confirming deterministic layer selection in practice.

\textbf{Stage 2 (h-e2 --- Quality Preservation Test):} Entropy-guided conversion of the 4 highest-entropy layers (indices 0, 1, 10, 31) to SWA(w = 512) produces a WikiText-103 perplexity of 6.3226 versus a full-attention baseline of 6.3386 --- a negligible decrease of 0.016 points, satisfying the $\leq$ 2.0-point gate criterion. This confirms that selective entropy-guided SWA conversion of 4 of 32 layers does not measurably degrade language modeling quality under the zero-shot condition.

These results establish that: (1) high-entropy layers in Llama-2-7B exist, can be stably identified from a small calibration set, and (2) converting the 4 highest-entropy layers to SWA(w = 512) preserves WikiText-103 perplexity within the target threshold. What remains to be determined is whether entropy-guided selection is superior to uninformed alternatives (random or last-k layer selection), and how quality degrades with more aggressive conversion ratios (k = 8). These questions are the direct next steps for follow-up experiments.

The broader implication is that pre-trained LLMs contain layers with already-diffuse attention patterns, and these layers can be identified cheaply at inference time without training data or gradient computation. Zero-shot selective conversion offers a path to reducing attention cost in deployed models without the overhead of fine-tuning.

